\chapter{HMC for State-Space Likelihoods}
\label{ch:bf-hmc-for-state-space}

This chapter develops fixed-trajectory Hamiltonian Monte Carlo for
filtering-based posterior targets.  The mathematical foundations follow the
standard HMC literature \citep{neal2011mcmc,betancourt2017conceptual}, but
implementation claims remain gated by BayesFilter's target, gradient, privacy,
and compilation contracts.

\section{Target Contract}
\label{sec:bf-hmc-target-contract}

For BayesFilter, an HMC transition is valid only relative to the target in
Definition~\ref{def:bf-filtering-target}.  The sampler needs the scalar log
target and its gradient in the coordinates used by the integrator.  If the
filter supplies a custom score, that score must correspond to the same log
likelihood value used in the Metropolis correction.

The HMC chapter therefore depends on the earlier derivative chapters.  A
filtering backend that can evaluate $\ell(\theta)$ but has no declared gradient
policy is value-only evidence.  It can be used for optimization, diagnostics,
or likelihood comparison, but it should not be promoted to HMC production use.

\section{Implementation Position}
\label{sec:bf-hmc-implementation-position}

TensorFlow Probability provides useful reference implementations of HMC and
NUTS, but BayesFilter does not treat TFP NUTS as the production tuning backend.
This is a project implementation decision, not a mathematical criticism of
NUTS \citep{hoffman2014nuts}.  The fixed-trajectory HMC path is the promoted
TensorFlow/TFP route for model-owned tuning because it keeps the value,
gradient, mass-matrix, step-size, and leapfrog-count contracts inspectable.

The reasons to keep TFP NUTS as a reference rather than the default production
path are:

\begin{itemize}
  \item full-trajectory NUTS control flow is hard to keep reliably XLA-compiled
    across model-specific likelihoods, adaptation logic, and diagnostics;
  \item nested TensorFlow Probability kernels can obscure whether the complete
    value-and-gradient pipeline is compiled or only partially traced;
  \item debugging numerical failures inside TFP kernel internals is slower than
    debugging a small BayesFilter-owned fixed-trajectory HMC transition;
  \item large DSGE targets require explicit policies for finite log density,
    finite gradients, static shapes, filter failure handling, and custom
    derivative backends.
\end{itemize}

Consequently, BayesFilter's production direction is a small, inspectable HMC
layer whose contracts are designed around filtering likelihoods,
analytic/custom gradients, and XLA/JIT discipline. Any future NUTS
implementation would need its own reviewed filter-aware contracts before it is
promoted beyond reference or diagnostic use.

\section{Promoted Fixed-Trajectory Tuning Algorithm}
\label{sec:bf-promoted-fixed-trajectory-hmc-tuning}

BayesFilter's promoted one-call tuning tool uses fixed-trajectory HMC rather
than a single one-dimensional search over the integration time.  The algorithm
is deliberately conservative:

\begin{enumerate}
  \item Build an initial mass matrix from posterior, pilot, or local covariance
    information.
  \item Regularize that covariance toward a symmetric positive-definite matrix;
    if dense regularization is not credible, use a recorded diagonal fallback.
  \item Choose an initial grid of leapfrog counts \(L\).
  \item For each \(L\), tune the step size \(\epsilon\) using bounded HMC
    screens and the target acceptance rule.
  \item Select the best pair \((L,\epsilon)\) using finite-value diagnostics,
    acceptance, trajectory-window membership, and hard-veto checks.
  \item If the selected \(L\) lies on a grid boundary, repair the grid upward or
    downward and repeat the screen.
  \item Run one final local grid around the selected pair.
  \item Freeze \((L,\epsilon)\), update the mass matrix moderately from a
    windowed draw, and repeat only while the repair loop records meaningful
    progress and an attempt slot remains.
\end{enumerate}

The step-size search and the trajectory length must refer to the same HMC
kernel.  In particular, a reasonable-step calculation performed with one
leapfrog step is not valid evidence for a later fixed trajectory with
\(L>1\).  BayesFilter therefore evaluates every reasonable-step proposal with
the active fixed \(L\), and records that value in the private attempt evidence.

Operational warmup also distinguishes a screened step from an unscreened
proposal.  If the bootstrap stage has already passed a fixed-kernel screen at
\((L,\epsilon_{\mathrm{screen}})\), the first warmup chart treats
\(\epsilon_{\mathrm{screen}}\) as a qualified upper bound.  Dual averaging may
adapt below that value, but the step actually supplied to the inner HMC kernel
and the retained handoff step are clamped so that
\[
  0 < \epsilon_t \leq \epsilon_{\mathrm{screen}}.
\]
The artifact records the raw dual-averaging proposal, the bounded next step,
the step consumed by HMC, and the number of ceiling hits.  Ceiling hits are
explanatory diagnostics: they show pressure toward an unqualified region, but
do not by themselves prove that the bound is optimal.

When a windowed mass update changes coordinates, the old step bound does not
transfer automatically.  BayesFilter maps the endpoint into the candidate
chart and qualifies a new fixed-\(L\) step using four independent momentum
probes per candidate.  All probes must execute with finite, target-valid
states; their acceptance probabilities are aggregated for the bracket.  A
model-domain failure, such as a non-positive-definite covariance at a proposed
parameter value, rejects that epsilon candidate and drives the bracket
downward.  An unrelated TensorFlow runner or device error remains an
implementation failure rather than being misclassified as low acceptance.

If a candidate chart cannot obtain a qualified step, the mass update is
rejected and the incumbent chart is retained.  This is a non-promoting metric
decision, not permission to weaken target-domain assertions or replace the
declared target.

This procedure is a tuning and diagnostics policy, not a posterior-validity
proof.  A passing tuner means that BayesFilter found a finite fixed HMC kernel
that survived its route-specific fresh verification.  For the ordinary default
runner, that handoff also requires the configured minimum retained draws and
finite rank-normalized split and folded split \(\widehat R\) values no greater
than 1.01.  The threshold is a tuning-handoff diagnostic; it does not establish
posterior convergence.  Chapter~\ref{ch:bf-hmc-tuning-interfaces} gives the
public interfaces and the exact distinction between tuning evidence and
retained-chain evidence.  Scientific validity, sampler superiority, default
readiness, and GPU readiness require separate evidence.

\section{Public And Private Handoffs}
\label{sec:bf-hmc-public-private-handoffs}

The final frozen kernel has two different artifacts.  The private BayesFilter
handoff contains the adapted mass arrays, step size, and leapfrog count needed
to replay the exact fixed kernel.  The public handoff is non-replayable: it may
expose hashes, signatures, target dimension, verification status, acceptance
summary, and nonclaims, but it must not expose raw samples, final states, mass
arrays, step size, leapfrog count, candidate grids, or private file paths.

This split lets downstream model repositories record that tuning completed
without forcing them to own or publish HMC mechanics.  If a retained frozen
kernel must be launched, the replay boundary remains BayesFilter-owned and
requires the private payload.

\section{Diagnostic Use of NUTS}
\label{sec:bf-nuts-diagnostic-use}

NUTS remains valuable as an algorithmic reference and as a diagnostic sampler on
small controlled targets.  The bounded BayesFilter decision is narrower: the
project should not treat TFP NUTS as the default answer to every convergence,
geometry, or compilation problem.  Larger filtering targets add the hard parts
omitted by small benchmarks: spectral derivative risk, missing-data branches,
model-specific invalid regions, and backend fallbacks.
